\paragraph{Support cellulaire du transport central.}
La récurrence précédente détermine aussi la géométrie finie de la
trajectoire. Nous considérons les positions immédiatement antérieures à chacune
des \(108\) étapes ; les coordonnées appartiennent à
\(D_9=\{1,\ldots,9\}\), avec un transport cyclique modulo neuf.

\begin{proposition}[Trois diagonales cycliques de la section centrale]
\label{prop:el-soporte-27}
La projection cellulaire de la trajectoire initiée en \((r,c,h,v)=(5,5,1,1)\) a pour support exact
\begin{equation}
 \begin{aligned}
 \operatorname{Supp}_{\mathrm{cell}}(\Gamma_e)&=\{(r,c)\in D_9^2:
                 c-r\equiv0,1,-1\pmod9\},
 \\
 |\operatorname{Supp}_{\mathrm{cell}}(\Gamma_e)|&=27.
 \end{aligned}
 \label{eq:el-soporte-27}
\end{equation}
L'orientation initiale conjuguée \((h,v)=(-1,-1)\), avec le même calendrier d'inversions, produit le même support.
\end{proposition}
\begin{proof}
Dans un bloc de \(27\) étapes, nous écrivons \(h=v=s\), où
\(s\in\{-1,1\}\). Si un bloc de trois étapes commence en \((a,a)\),
la lecture additive précède le déplacement jusqu'à \((a,a+s)\) ;
la lecture multiplicative précède le déplacement jusqu'à
\((a+s,a+s)\) ; le seuil conserve cette dernière position. Les trois
états antérieurs appartiennent donc aux diagonales
\(c-r=0\) et \(c-r=s\) du tore positionnel. Les neuf blocs
consécutifs parcourent les neuf possibilités de \(a\), car la
translation par \(s\) est d'ordre neuf. L'inversion suivante
parcourt la diagonale adjacente de signe contraire. Les diagonales
de différences \(0,1,-1\) sont disjointes, chacune contient neuf
cellules et toutes sont visitées. Inverser le signe initial
inverse l'ordre de visite des deux diagonales adjacentes.
\end{proof}

Le point fixe \((5,5)\) caractérise la section d'origine ; l'ensemble \(\operatorname{Supp}_{\mathrm{cell}}(\Gamma_e)\) caractérise sa trajectoire cellulaire. Le résultat détermine une morphologie discrète sur le domaine APP. La réalisation spatiale et le recollement d'orientation reçoivent en outre les données de transport qui leur correspondent. La permanence de la section et la variation de sa position observable sont ainsi des propriétés compatibles.

\paragraph{Deux historiques centraux dans la coordinatisation intégrale.}
La coordinatisation \(\mathcal L_{12}\) de la
section~\ref{subsec:k-registro-integral} conserve les événements et leur
ordre. Nous l'appliquons aux deux orientations centrales. Pour chaque
pas \(t\), soit \(d_t=\rho_9(u_t)\), où \(u_t\) est l'entier
additif, multiplicatif ou de seuil lu par la récurrence.
Dans la fenêtre \(E_{i+1}=\{9i+1,\ldots,9i+9\}\), nous définissons
\[
 n_i=\#\{t\in E_{i+1}:d_t\in\{2,4,6,8\}\},\qquad
 e_i=\Sigma_{12}(i+1)n_i,\qquad0\le i<12.
\]
Le signe \(\Sigma_{12}\) est l'orientation sectorielle de la coordinatisation.
Dans la trajectoire \({++}\), il coïncide avec le signe horizontal
transporté ; dans la trajectoire \({--}\), il coïncide avec son opposé.
De manière équivalente, dans les deux trajectoires, il est
\(\Sigma_{12}(i+1)=h^{\mathrm{pre}}_{9i+1}/h_0\), où
\(h^{\mathrm{pre}}_t\) est l'orientation antérieure à la lecture du pas.
Cette normalisation conserve séparément l'orientation initiale et
l'orientation relative. Les symboles \(e_i\) désignent des événements
entiers, tandis que \(e\) sans indice désigne le nombre d'Euler.

L'évaluation par fenêtres des deux récurrences est
\[
\begin{array}{c|rrrrrrrrrrrr}
i&0&1&2&3&4&5&6&7&8&9&10&11\\\hline
n_i^{++}&2&2&5&4&3&2&2&2&5&4&3&2\\
n_i^{--}&4&3&2&2&2&5&4&3&2&2&2&5
\end{array}
\]
et détermine les histoires
\begin{equation}
 e_{++}=(2,2,5,-4,-3,-2)^2,\qquad
 e_{--}=(4,3,2,-2,-2,-5)^2.
 \label{eq:el-historias-centrales}
\end{equation}
La puissance indique la concaténation. Les deux histoires ont une somme nulle,
la même suite de signes sectoriels et le même histogramme
en valeur absolue : six deux, deux trois, deux quatre et deux cinq.

\begin{proposition}[Séparation réversible des deux histoires]
\label{prop:el-memorias-distintas}
Dans les coordonnées \((s_C,M,a)=\mathcal L_{12}(e)\), les deux histoires
satisfont \(s_C=0\) et \(a=0\), tandis que
\begin{equation}
 M_{++}=(8,8,14,-4,-2),\qquad
 M_{--}=(18,16,14,6,6).
 \label{eq:el-memorias-distintas}
\end{equation}
La coordinatisation reconstruit chaque historique de manière unique.
\end{proposition}
\begin{proof}
Les deux sextuplets se répètent, de sorte que
\(p_i=e_i+e_{i+6}=2e_i\) et \(a_i=e_i-e_{i+6}=0\).
Soustraire \(p_5\) aux cinq premiers \(p_i\) produit les marges
de \eqref{eq:el-memorias-distintas}. Leurs sommes sont \(24\) et
\(60\), respectivement ; la formule inverse du
lemme~\ref{lem:k-carta-integral} donne \(p_5=-4\) et \(p_5=-10\).
Ensuite, \(p_i=p_5+M_i\) et
\(e_i=e_{i+6}=p_i/2\) restituent exactement les deux sextuplets.
La divisibilité par six et les six congruences de
parité qui caractérisent l'image entière du lecteur sont satisfaites.
\end{proof}

L'égalité des recensements et des signes conserve une projection commune ; les
marges conservent la différence d'historique. La récurrence opposée
émet \(\gamma_{--}=(b^3a^3)^2\) : après \(27\) pas, la route
\({++}\) atteint de nouveau le centre avec l'orientation initiale de
\({--}\), et le calendrier de trois pas conserve sa phase. Ainsi,
\(\gamma_{--,t}=\gamma_{e,t+27}\), avec des indices cycliques.
Les dénombrements \(D_k\) sont invariants sous ce décalage ; le
lecteur de fenêtres reçoit le même mot sectoriel dans les deux
historiques. Les deux lecteurs de retour coïncident donc pour
ces registres que la coordinatisation intégrale distingue. L'évaluation massique et
la reconstruction des événements possèdent des portées informationnelles
précises au sein du même transport.
